\documentclass[12pt,a4paper]{article}

\usepackage[utf8]{inputenc}
\usepackage{geometry}
\geometry{margin=1in}
\usepackage{fancyhdr}
\usepackage{graphicx}
\usepackage{listings}
\usepackage{xcolor}
\usepackage{hyperref}
\usepackage{booktabs}
\usepackage{amsmath}

\hypersetup{
    colorlinks=true,
    linkcolor=blue,
    filecolor=magenta,
    urlcolor=cyan,
}

\lstset{
    backgroundcolor=\color{gray!10},
    basicstyle=\ttfamily\small,
    breaklines=true,
    frame=single,
    captionpos=b,
    numbers=left,
    numberstyle=\tiny\color{gray},
}

\pagestyle{fancy}
\fancyhf{}
\rhead{\thepage}
\lhead{SmartStudy - Project Report}

\begin{document}

\begin{titlepage}
    \centering
    \vspace*{2cm}

    {\LARGE\bfseries SmartStudy\par}
    \vspace{0.5cm}
    {\large Study Session Manager\par}
    \vspace{2cm}

    {\Large\bfseries Project Report\par}
    \vspace{2cm}

    {\large Bring Your Own Project (BYOP)\par}
    \vspace{1cm}

    \textbf{Student Details}\par
    \vspace{0.5cm}
    {\large Name: Shaikh Mohammad Warsi\par}
    {\large Registration No.: 24BAI10046\par}
    {\large University: VIT Bhopal\par}
    {\large Course: Vityarthi\par}
    \vspace{2cm}

    \today

    \vspace{3cm}

    \begin{abstract}
        This report presents SmartStudy, a JavaFX-based desktop application designed to help students manage their study sessions using the Pomodoro technique. The application provides a customizable study timer, tracks completed sessions, and offers analytics on study habits with personalized recommendations.
    \end{abstract}

\end{titlepage}

\tableofcontents
\newpage

\section{1. Introduction}

In modern academic life, students often struggle with maintaining consistent study habits and managing their time effectively across multiple subjects. Traditional study methods lack structure and fail to provide feedback on actual study progress. Without a proper system, students may overestimate or underestimate their study time, leading to poor planning and reduced productivity.

\textbf{SmartStudy} is a JavaFX-based desktop application designed to help students manage their study sessions using the proven Pomodoro technique. The application provides a customizable study timer, tracks completed sessions, and offers analytics on study habits. By recording and analyzing study patterns, SmartStudy helps students understand how they spend their time and provides personalized recommendations for improvement.

This project was developed as part of the Bring Your Own Project (BYOP) college course requirement. It demonstrates practical application of Java programming, Object-Oriented Programming (OOP) principles, file handling, and Graphical User Interface (GUI) development using JavaFX.

\section{2. Problem Identification}

Through personal experience and observation of fellow students, several key problems were identified in the context of self-directed study:

\subsection{2.1 Lack of Structured Study Time}
Many students study without a clear plan or time structure. Without defined study sessions, time can slip away without meaningful progress. The Pomodoro technique addresses this by breaking study time into focused intervals, but implementing it manually is inconvenient.

\subsection{2.2 No Progress Visibility}
Students rarely have a clear picture of how much time they actually spend studying each subject. This makes it difficult to allocate time effectively or identify imbalanced study habits.

\subsection{2.3 No Accountability System}
Without tracking, students cannot review their past study sessions or measure improvement over time. This lack of data makes it hard to identify problems or celebrate progress.

\subsection{2.4 Difficulty Maintaining Breaks}
Even when students do take breaks, they often extend too long. A system that suggests appropriate break durations based on study time can help maintain balance.

\section{3. Objectives}

The primary objectives of SmartStudy are:

\subsection{3.1 Core Functionality}
\begin{itemize}
    \item Develop a functional desktop application with an intuitive user interface
    \item Implement a customizable Pomodoro-based study timer (5-120 minutes)
    \item Create a session tracking system that records subject, duration, and notes
    \item Provide data persistence using local JSON file storage
\end{itemize}

\subsection{3.2 Analytics Features}
\begin{itemize}
    \item Calculate and display statistics (daily, weekly, monthly study time)
    \item Show subject-wise time breakdown
    \item Generate bar charts for visual progress tracking
    \item Provide personalized study recommendations
\end{itemize}

\subsection{3.3 Technical Objectives}
\begin{itemize}
    \item Demonstrate understanding of Java OOP concepts
    \item Implement proper MVC architecture with service layer
    \item Use JavaFX for GUI development
    \item Handle file I/O for data persistence
\end{itemize}

\section{4. Proposed Solution}

\textbf{SmartStudy} proposes a desktop application that combines:

\subsection{4.1 Smart Timer}
A fully controllable timer with start, pause, and reset functionality. Users can customize duration (5-120 minutes) and select the subject they're studying.

\subsection{4.2 Session Tracking}
Each completed session is automatically saved with subject, duration, start time, end time, and optional notes.

\subsection{4.3 Statistics Dashboard}
A dedicated statistics tab showing:
\begin{itemize}
    \item Today's, this week's, and total study time
    \item Subject breakdown chart
    \item Daily study time bar chart
    \item AI-generated recommendations
\end{itemize}

\subsection{4.4 Data Persistence}
All session data is stored locally in JSON format, allowing data to persist between application runs without requiring a database server.

\section{5. System Design}

\subsection{5.1 Architecture Overview}

SmartStudy follows a layered architecture pattern as illustrated in Figure~\ref{fig:architecture}.

\begin{figure}[htbp]
    \centering
    \begin{tabular}{|c|}
        \hline
        \textbf{Presentation Layer} \\
        \hline
        \textbf{Service Layer} \\
        \hline
        \textbf{Data Access Layer} \\
        \hline
        \textbf{Model Layer} \\
        \hline
    \end{tabular}
    \caption{System Architecture Layers}
    \label{fig:architecture}
\end{figure}

\subsection{5.2 Module Description}

\textbf{Models (\texttt{com.smartstudy.models})}
\begin{itemize}
    \item \texttt{StudySession}: Represents a single study session with subject, duration, notes, timestamps
    \item \texttt{Subject}: Represents a study subject with name, color, and aggregated statistics
    \item \texttt{StudyStatistics}: Container for calculated statistics
\end{itemize}

\textbf{Services (\texttt{com.smartstudy.services})}
\begin{itemize}
    \item \texttt{SessionService}: Handles CRUD operations for study sessions
    \item \texttt{AnalyticsService}: Calculates statistics and generates recommendations
\end{itemize}

\textbf{Utilities (\texttt{com.smartstudy.utils})}
\begin{itemize}
    \item \texttt{FileManager}: Handles JSON file read/write operations
    \item \texttt{DateUtils}: Provides date formatting and calculation utilities
\end{itemize}

\textbf{UI Layer (\texttt{com.smartstudy.ui})}
\begin{itemize}
    \item \texttt{SmartStudyApplication}: JavaFX application entry point
    \item \texttt{MainController}: Handles all UI interactions and updates
\end{itemize}

\subsection{5.3 Data Flow}

\begin{enumerate}
    \item User selects subject and duration on the Study Session tab
    \item User clicks Start, creating a new StudySession object
    \item Timer counts down using JavaFX Timeline animation
    \item On completion, session is saved via SessionService
    \item SessionService persists data to JSON via FileManager
    \item Statistics tab queries SessionService, AnalyticsService calculates results
    \item Charts and recommendations are displayed
\end{enumerate}

\subsection{5.4 Project Structure}

\begin{verbatim}
smartstudy/
|-- src/main/java/com/smartstudy/
|   |-- SmartStudyApplication.java
|   |-- models/
|   |   |-- StudySession.java
|   |   |-- Subject.java
|   |   `-- StudyStatistics.java
|   |-- services/
|   |   |-- SessionService.java
|   |   `-- AnalyticsService.java
|   |-- utils/
|   |   |-- FileManager.java
|   |   `-- DateUtils.java
|   `-- ui/
|       `-- controllers/
|           `-- MainController.java
|-- src/main/resources/
|   |-- fxml/
|   |   `-- main.fxml
|   `-- css/
|       `-- styles.css
|-- data/
|-- docs/
|-- pom.xml
`-- README.md
\end{verbatim}

\section{6. Implementation Details}

\subsection{6.1 Timer Implementation}

The timer uses JavaFX's \texttt{Timeline} class with a \texttt{KeyFrame} that fires every second:

\begin{lstlisting}[language=Java, caption={Timer Implementation}]
timer = new Timeline(new KeyFrame(Duration.seconds(1), e -> {
    if (remainingSeconds > 0) {
        remainingSeconds--;
        timerLabel.setText(DateUtils.formatTimerDisplay(remainingSeconds));
    } else {
        timerComplete();
    }
}));
timer.setCycleCount(Animation.INDEFINITE);
\end{lstlisting}

\subsection{6.2 Data Persistence}

Sessions are stored as JSON using Gson library:

\begin{lstlisting}[language=Java, caption={JSON Persistence}]
public static void saveSessions(List<StudySession> sessions) {
    ensureDataDirectoryExists();
    try (Writer writer = new FileWriter(SESSIONS_FILE)) {
        gson.toJson(sessions, writer);
    } catch (IOException e) {
        System.err.println("Error saving sessions: " + e.getMessage());
    }
}
\end{lstlisting}

\subsection{6.3 Statistics Calculation}

The AnalyticsService aggregates session data:

\begin{lstlisting}[language=Java, caption={Statistics Calculation}]
public StudyStatistics calculateStatistics() {
    StudyStatistics stats = new StudyStatistics();
    stats.setTotalSessions(sessions.stream()
            .filter(StudySession::isCompleted)
            .count());
    stats.setTodayMinutes(getMinutesForPeriod("today"));
    // ... additional calculations
    return stats;
}
\end{lstlisting}

\subsection{6.4 Recommendations Engine}

Simple rule-based recommendations:

\begin{lstlisting}[language=Java, caption={Recommendations Engine}]
public List<String> getRecommendations() {
    if (sessions.size() < 5) {
        return List.of("Start with 2-3 short study sessions daily.");
    }
    // Analyze patterns and generate suggestions
}
\end{lstlisting}

\section{7. Challenges Faced}

\subsection{7.1 Technical Challenges}

\textbf{JavaFX Setup}: Setting up JavaFX with Maven required correct dependency configuration and module-path arguments. The pom.xml had to be configured properly with the JavaFX Maven plugin.

\textbf{Timer Accuracy}: Using \texttt{Timeline} for the timer required careful handling to ensure the display updates correctly and the timer completes at the right moment.

\textbf{JSON Serialization}: The \texttt{LocalDateTime} fields required custom handling since Gson doesn't serialize them by default. The model classes needed proper structure.

\subsection{7.2 Design Challenges}

\textbf{Balancing Features}: Deciding which features to include while keeping the project achievable within the timeframe required careful prioritization.

\textbf{User Experience}: Creating an intuitive interface without overwhelming users with options was challenging. The tab-based layout helps organize functionality.

\subsection{7.3 Learning Challenges}

\textbf{JavaFX FXML}: Learning the FXML syntax and understanding how it connects to controller classes required careful study of the framework.

\textbf{MVC Pattern}: Properly separating concerns between the UI layer and business logic was important but sometimes challenging to get right.

\section{8. Results}

\subsection{8.1 Functional Results}

SmartStudy successfully implements all planned features:

\begin{itemize}
    \item[$\checkmark$] Customizable Pomodoro timer (5-120 minutes)
    \item[$\checkmark$] Subject selection with 9 predefined options
    \item[$\checkmark$] Start, pause, and reset timer controls
    \item[$\checkmark$] Session notes functionality
    \item[$\checkmark$] Automatic session saving on completion
    \item[$\checkmark$] JSON-based data persistence
    \item[$\checkmark$] Daily, weekly, and monthly statistics
    \item[$\checkmark$] Subject breakdown display
    \item[$\checkmark$] Bar chart visualization
    \item[$\checkmark$] Personalized recommendations
\end{itemize}

\subsection{8.2 Technical Validation}

\begin{itemize}
    \item Application compiles without errors using Maven
    \item All MVC layers properly separated
    \item Data persists correctly between runs
    \item JavaFX GUI renders as designed
    \item Charts display data correctly
\end{itemize}

\subsection{8.3 User Experience}

\begin{itemize}
    \item Clean, intuitive tab-based interface
    \item Responsive timer controls
    \item Clear visual feedback during sessions
    \item Informative statistics dashboard
    \item Helpful recommendations
\end{itemize}

\section{9. Future Scope}

While SmartStudy addresses the core needs of study session management, several enhancements could be considered for future versions:

\subsection{9.1 Extended Features}
\begin{itemize}
    \item Export statistics to PDF or Excel
    \item Add custom subjects with color coding
    \item Set daily study goals with notifications
    \item Track assignment deadlines
\end{itemize}

\subsection{9.2 Data Enhancements}
\begin{itemize}
    \item Cloud synchronization across devices
    \item Backup and restore functionality
    \item Data import/export options
\end{itemize}

\subsection{9.3 Advanced Analytics}
\begin{itemize}
    \item Weekly study pattern analysis
    \item Performance prediction based on study time
    \item Spaced repetition scheduling integration
\end{itemize}

\subsection{9.4 Platform Expansion}
\begin{itemize}
    \item Mobile companion app (Android/iOS)
    \item Web version with cloud backend
    \item Browser extension for web-based study resources
\end{itemize}

\section{10. Conclusion}

SmartStudy represents a practical application of Java programming concepts learned throughout the course. By implementing a real-world tool for study management, the project demonstrates proficiency in several key areas:

\begin{enumerate}
    \item \textbf{Java Programming}: Core language features including collections, streams, and lambda expressions
    \item \textbf{Object-Oriented Design}: Proper use of classes, encapsulation, and separation of concerns
    \item \textbf{JavaFX GUI Development}: Building user interfaces with FXML and styling with CSS
    \item \textbf{File Handling}: JSON-based data persistence using Gson
    \item \textbf{Software Architecture}: MVC pattern with service layer
\end{enumerate}

The project successfully meets the BYOP course requirements while creating a genuinely useful tool that students can employ in their academic journey. The Pomodoro timer and progress tracking features provide real value, while the analytics and recommendations demonstrate practical application of algorithmic thinking.

\section{References}

\begin{enumerate}
    \item JavaFX Documentation - \url{https://openjfx.io/}
    \item Gson Library - \url{https://github.com/google/gson}
    \item Maven Documentation - \url{https://maven.apache.org/}
    \item The Pomodoro Technique - Francesco Cirillo
    \item Java 17 API Documentation - \url{https://docs.oracle.com/en/java/javase/17/}
\end{enumerate}

\vfill

\noindent\textit{Project submitted as part of BYOP course at VIT Bhopal} \\
\textit{Student: Shaikh Mohammad Warsi (24BAI10046)} \\
\textit{GitHub Repository: https://github.com/yourusername/smartstudy}

\end{document}
